\documentclass[11pt]{article}
\usepackage[margin=1in]{geometry}
\usepackage{amsmath,amssymb,amsthm,mathtools}
\usepackage{microtype}
\usepackage{enumitem}
\usepackage{xcolor}
\usepackage{hyperref}
\hypersetup{colorlinks=true,linkcolor=blue!55!black,citecolor=blue!55!black}

% ---- certainty-tier tags (consistent across the DBP corpus) ----
\newcommand{\Proven}{\textsf{[proven]}}
\newcommand{\Symbolic}{\textsf{[symbolic]}}
\newcommand{\Measured}{\textsf{[computer-verified]}}
\newcommand{\Conj}{\textsf{[conjecture]}}
\newcommand{\Interp}{\textsf{[interpretation]}}

\theoremstyle{plain}
\newtheorem{theorem}{Theorem}
\newtheorem{corollary}{Corollary}
\newtheorem{proposition}{Proposition}
\theoremstyle{definition}
\newtheorem{definition}{Definition}
\newtheorem{example}{Example}

% ---- standardised notation (see Notation note) ----
\newcommand{\KG}{K_{G}}
\newcommand{\Ac}{\mathcal{A}_c}
\newcommand{\kc}{\kappa_{c}}
\newcommand{\kint}{\kappa_{\mathrm{int}}}
\newcommand{\ks}{\kappa_{s}}
\newcommand{\kcrole}[1]{\kappa_{c}^{(#1)}}   % chart-indexed coupling channel
\newcommand{\Reg}{\Omega}
\newcommand{\jet}{j}

\title{\textbf{The Active Role-Jet Orbit Calculus for Constraint-Surface Roles}}
\author{William Lloyd}
\date{\today}

\begin{document}
\maketitle

\begin{abstract}
\noindent This paper develops the foundational calculus of \emph{Directive--Substrate--Product}
(DBP) role data on a local three-role relation $F(D,S,P)=0$. The central object is not a passive
permutation of coordinate labels but an \emph{active} family of charts: the same relation re-solved
with each coordinate as output, with finite jet data transported through the resulting
inverse-function maps. On the regular locus this active transposition is an exact birational
representation of $S_3$ on finite jets, localised at the role denominators. The \emph{Active
Role-Jet Orbit Theorem} identifies the resulting orbit as the canonical carrier of local DBP
invariants: a functional is a DBP invariant exactly when it factors through the orbit quotient,
so invariants are obtained, not postulated. The same construction yields a channel-reduction exact
sequence in which a channel account is a gauge- and frame-relative representative over an intrinsic
invariant quotient; in the second-order curvature case this is the three-channel Gaussian
decomposition $\KG=\kc+\kint+\ks$. From the explicit second-jet transport we obtain the named
output-role coupling channels $\Lambda_P,\Lambda_D,\Lambda_S$, the coupling carrier
$O=(\kcrole{P},\kcrole{D},\kcrole{S})$, and the role anisotropy scalar $\Ac$. A faithfulness
corollary shows the carrier map has Jacobian determinant $8\,\Lambda_P\Lambda_D\Lambda_S/q_0^6$,
so the carrier separates active role orbits off the channel-isotropy loci. Singular points where an
output chart cannot be solved are treated as typed strata, not removable cases. Every claim is
tagged by certainty tier; the construction is self-contained.
\end{abstract}

\paragraph{Notation.} Throughout, $\KG$ is the Gaussian curvature. For a graph $P=f(D,S)$ we write
the first and second jet as $(a,b,A,B,C)=(f_D,f_S,f_{DD},f_{DS},f_{SS})$ and set
$q_0=1+a^2+b^2$; for the implicit relation $F=P-f$ one has $g=\nabla F$ and $q_0=g^{\mathsf T}g$,
so the graph normaliser and the implicit gradient norm coincide in the graph gauge. We distinguish
two ``three-channel'' structures and never conflate them: the channel \emph{types}
$\kc,\kint,\ks$ (coupling, interaction, self) decompose $\KG$ \emph{within one chart}; the
chart-indexed coupling channels $\kcrole{\rho}$, $\rho\in\{P,D,S\}$, are the coupling type evaluated
in the three output-role \emph{charts}, and the anisotropy $\Ac$ is their spread. ``The three
channels'' is never used unqualified.

%==================================================================
\section{Introduction: active versus passive role data}
\label{sec:intro}

Let $(D,S,P)$ denote a distinguished ordered triple of local coordinates. DBP studies a local
relation among them by comparing the finite jets obtained when the same relation is solved with
different coordinates as output. The basic object is therefore neither a labelled triple alone nor a
passive permutation of a fixed derivative array; it is an active family of charts related by
inverse-function maps.

This distinction is essential. If one merely permutes coordinate labels, a symmetric expression is
unchanged for the trivial reason that the indices were renamed; the calculation verifies covariance
but measures no role dependence. Active role transposition is different: it changes the graph
representative of the same local relation. The derivative coordinates then transform rationally and
need not agree componentwise.

The development below is mechanical. First, the relation is restricted to the regular locus where
each output chart exists (\S\ref{sec:roledata}). Second, finite jets are transported by the active
role maps, which realise an exact birational $S_3$ action (\S\ref{sec:transpose}). Third, the active
orbit is used as the carrier for local DBP invariants (\S\ref{sec:orbit}); the quotient property of
invariants is standard, but the explicit active action, its rational finite-jet closure, its
separation from passive covariance, and the induced channel reduction are the content. The remaining
sections specialise the channel reduction to curvature (\S\ref{sec:channel}), extract the named
output-role channels and prove faithfulness of the coupling carrier (\S\ref{sec:channels-named}),
record gauge transport (\S\ref{sec:transport}), type the singular strata (\S\ref{sec:singular}), and
delimit scope (\S\ref{sec:scope}).

%==================================================================
\section{Local role data}
\label{sec:roledata}

\subsection{Three-role relation}
Let $U\subset\mathbb R^3$ be open with coordinates $(D,S,P)$, let $F\in C^{r+1}(U)$, and let
$x_0=(D_0,S_0,P_0)\in U$ satisfy $F(x_0)=0$. A local graph chart with $P$ as output is written
$P=f(D,S)$. The implicit and graph descriptions agree on any chart in which the relevant partial of
$F$ is nonzero. We write $\jet^r f$ for the order-$r$ jet of the graph representative and $\jet^rF$
for the order-$r$ implicit jet.

\subsection{Regular active-role locus}
Active role transposition requires the relation to be solvable with each coordinate as output. Define
the first-order regular role locus
\[
\Reg_1=\{x\in U: F(x)=0,\ F_D(x)F_S(x)F_P(x)\ne0\},
\]
and more generally let $\Reg_r$ denote the locus on which all output charts needed for the order-$r$
active role object exist. In graph form $P=f(D,S)$ with $a=f_D,\ b=f_S$, full first-order active
recharting is defined where
\[
a\ne0,\qquad b\ne0 .
\]
These are the denominators of the inverse-function maps, not technical decorations: a point outside
the active-role locus is a singular stratum for the requested role object (\S\ref{sec:singular}).

%==================================================================
\section{Active role transposition: the birational \texorpdfstring{$S_3$}{S3}}
\label{sec:transpose}

\subsection{First-order action}
For a graph chart $P=f(D,S)$ with $a=f_D,\ b=f_S$, the input swap $s=(D\;S)$ and the active
transposition $t=(D\;P)$ act by
\[
s(a,b)=(b,a),
\qquad
t(a,b)=\Bigl(\tfrac1a,\,-\tfrac ba\Bigr),
\]
obtained by re-solving $F=0$ for $D$ as output. These satisfy
\[
s^2=t^2=e,\qquad (st)^3=e,
\]
so active role transposition is an exact birational representation of $S_3$ on the regular jet locus,
localised at the role denominators. \Proven

\subsection{Second-order generator}
With $A=f_{DD},\ B=f_{DS},\ C=f_{SS}$, the input swap acts by
\[
s(a,b,A,B,C)=(b,a,C,B,A).
\]
Under $t=(D\;P)$ the chart $D=q(P,S)$ has
\[
q_P=\frac1a,\qquad q_S=-\frac ba,\qquad
q_{PP}=-\frac{A}{a^3},\qquad
q_{PS}=\frac{Ab-aB}{a^3},\qquad
q_{SS}=\frac{-Ab^2+2abB-a^2C}{a^3},
\]
so that
\[
t(a,b,A,B,C)=
\left(\frac1a,\,-\frac ba,\,-\frac{A}{a^3},\,\frac{Ab-aB}{a^3},\,
\frac{-Ab^2+2abB-a^2C}{a^3}\right).
\]
The maps again satisfy the $S_3$ relations on the regular locus. The mixed second-order numerator
$Ab-aB$ appearing in $q_{PS}$ is the source of the directive coupling channel of
\S\ref{sec:channels-named}. Higher-order generators follow by the finite inverse-function theorem and
the chain rule; at each finite order they are rational in the original jet, localised at the
first-order role denominators (Appendix~\ref{app:higher}). \Proven

%==================================================================
\section{The Active Role-Jet Orbit Theorem}
\label{sec:orbit}

\begin{definition}[jet carrier]
Let $J_r(x_0)$ be the order-$r$ jet of the chosen local graph representative at $x_0$. The active
order-$r$ role-jet orbit is $\mathfrak R_r(x_0)=G\cdot J_r(x_0)$, where the action is active role
recharting. Retaining the coordinate values gives the value-plus-jet orbit
$\mathfrak R_r^{+}(x_0)=G\cdot(x_0,J_r(x_0))$.
\end{definition}

\begin{definition}[local DBP invariant]
Let $X_r$ be the order-$r$ jet space on the regular role locus and $Y$ any target set. A functional
$\phi:X_r\to Y$ is an order-$r$ DBP invariant if $\phi(\mathcal T_\sigma\xi)=\phi(\xi)$ for all
$\sigma\in S_3$ wherever both sides are defined. For value-plus-jet invariants, replace $X_r$ by
$X_r^{+}$.
\end{definition}

\begin{theorem}[Active Role-Jet Orbit Theorem]
\label{thm:orbit}
On the regular active-role locus, let $G=S_3$ act on $X_r$ by the active finite-jet recharting maps
of \S\ref{sec:transpose}, and let $\pi:X_r\to X_r/G$ be the orbit quotient. Then $\phi:X_r\to Y$ is a
local order-$r$ DBP invariant if and only if there is a unique $\psi:X_r/G\to Y$ with
$\phi=\psi\circ\pi$. Equivalently, the value of any order-$r$ DBP invariant at $x_0$ is a function of
the active orbit $\mathfrak R_r(x_0)=G\cdot J_r(x_0)$. The same holds for value-plus-jet functionals
with $X_r^{+}$ and $\mathfrak R_r^{+}(x_0)$. \Proven
\end{theorem}

\begin{proof}
The maps $\mathcal T_\sigma$ define an action of $S_3$ on the regular active-role locus. By
definition a DBP invariant is constant on its orbits. The quotient map $\pi$ identifies exactly the
elements lying in the same orbit; hence a functional is constant on active role orbits iff it factors
uniquely through the quotient. The value-plus-jet statement is identical with $X_r^{+}$. The quotient
step is formal; the mathematical content is the prior construction of the active finite-jet action,
in which the representatives in different output roles are rational inverse-function transforms, not
passive relabellings.
\end{proof}

%==================================================================
\section{Consequences of the orbit carrier}
\label{sec:consequences}

\subsection{Completeness criterion}
A finite tuple $K_r$ is a complete order-$r$ DBP kernel only if it separates the same active role
orbits as the quotient map:
\[
K_r(\xi_1)=K_r(\xi_2)\quad\Longleftrightarrow\quad \pi(\xi_1)=\pi(\xi_2)
\]
on the intended domain. A preferred kernel cannot be assumed complete merely because its entries are
invariant; completeness requires orbit separation. \Proven\ This criterion is the test applied to the
coupling carrier in \S\ref{sec:faithful}.

\subsection{Exact rational closure}
If the original finite jet has entries in a field $\mathbb K$ and all active role denominators are
nonzero, every active role-transposed jet has entries in the localisation obtained by inverting those
denominators. In particular
\[
J_r(x_0)\in\mathbb Q\quad\Longrightarrow\quad G\cdot J_r(x_0)\subset\mathbb Q[\Delta^{-1}],
\]
where $\Delta$ is the product of the role denominators used by the active charts. \Proven\ Thus for
exact rational input jets, active role transport is an exact rational operation on the regular locus;
this is the statement that licenses exact-$\mathbb Q$ computation of every quantity below.

\subsection{Descent classes}
Invariants separate by the orbit data they require. \emph{Value-only} invariants descend to the
quotient of coordinate values alone; \emph{jet-only} coupling invariants descend to the quotient of
the jet without values; \emph{mixed} value-jet invariants require the value-plus-jet orbit
$\mathfrak R_r^{+}$. The curvature channels of \S\ref{sec:channel} are jet-only; the power-law degree
spectra of Appendix~\ref{app:examples} are value-only in log-role form. \Proven

%==================================================================
\section{Channel reduction}
\label{sec:channel}

\subsection{Abstract reduction law}
Let $C_r$ be a channel-account space at order $r$ and let $\Sigma:C_r\to I_r$ be a surjective
reduction onto an invariant range. Then
\[
0\longrightarrow\ker\Sigma\longrightarrow C_r\xrightarrow{\ \Sigma\ }I_r\longrightarrow0
\]
is exact, and $C_r/\ker\Sigma\cong I_r$ is the invariant. The fibre over a fixed value of $I_r$ is the
channel account: a representative may move inside the fibre under chart or gauge change while the
reduced invariant is fixed. \Proven

\subsection{Curvature-channel projection}
For a hypersurface relation $F=0$ at a regular point set $g=\nabla F$, $H=\nabla^2F$, $q=g^{\mathsf
T}g\ne0$. In the distinguished coordinate basis split
\[
H=H_c+H_s,\qquad H_s=\operatorname{diag}(H),\qquad H_c=H-H_s .
\]
For $1\le r\le n-1$ define the channel-density polynomial
\[
\widehat{\mathcal C}_r(t,u)
=(-1)^{r+1}\sum_{|I|=r+1}\det
\begin{pmatrix}0 & g_I^{\mathsf T}\\ g_I & (tH_c+uH_s)_I\end{pmatrix}
=\sum_{p+q=r}t^p u^q\,\widehat\kappa_{r;p,q}.
\]
The normalised channel curvatures are $\kappa_{r;p,q}=\widehat\kappa_{r;p,q}/q^{(r+2)/2}$, and the
reduced elementary curvature invariant is
\[
\sigma_r=\sum_{p+q=r}\kappa_{r;p,q}=\frac{\widehat{\mathcal C}_r(1,1)}{q^{(r+2)/2}} .
\]
Thus the reduction is summation, $\Sigma:C_r\mapsto\sum_{p+q=r}\kappa_{r;p,q}=\sigma_r$: the channel
vector is the representative, its sum is the invariant. \Proven

\subsection{The three-variable Gaussian case}
For $n=3$, $r=2$ the sum has three terms; writing $\kc=\kappa_{2;2,0}$ (the $t^2$, pure-coupling
term), $\kint=\kappa_{2;1,1}$ (the $tu$ cross term), $\ks=\kappa_{2;0,2}$ (the $u^2$, pure-self term),
\[
C_2=(\kc,\kint,\ks),\qquad \KG=\kc+\kint+\ks .
\]
The zero-sum transport plane is $\ker\Sigma=\{(u,v,w):u+v+w=0\}$: a gauge or chart change may move the
channel triple within the affine plane over $\KG$, but cannot change the reduced Gaussian curvature.
\Proven

\subsection{Graph-channel formula}
For any local graph $z=h(x,y)$ with $\alpha=h_x,\ \beta=h_y$, $L=h_{xx},\ M=h_{xy},\ N=h_{yy}$ and
$Q=1+\alpha^2+\beta^2$,
\[
\kc=-\frac{M^2}{Q^2},\qquad \ks=\frac{LN}{Q^2},\qquad \kint=0,\qquad
\KG=\frac{LN-M^2}{Q^2}.
\]
Every graph-normalised output chart therefore has zero graph-channel interaction; interaction
reappears in non-graph implicit gauges or under gauge transport between graph sections
(\S\ref{sec:transport}). \Proven

\subsection{Ambient-dimension channel grid}
The bordered-minor law is already an ambient-$n$ statement, not merely a surface identity.  At
$r=2$, let $\widehat{\mathcal C}^{\,I}_2(t,u)$ denote its summand for a coordinate triple
$I=\{i,j,k\}$.  Coefficient extraction gives the local-to-ambient triple
\[
\bigl(\widehat\kappa_{2;2,0},\widehat\kappa_{2;1,1},\widehat\kappa_{2;0,2}\bigr)
=\sum_{i<j<k}
\bigl([t^2]\widehat{\mathcal C}^{\,ijk}_2,
      [tu]\widehat{\mathcal C}^{\,ijk}_2,
      [u^2]\widehat{\mathcal C}^{\,ijk}_2\bigr),
\]
and division by $q^2$ gives the corresponding normalised channels.  Hence the familiar
$n=3$ Gaussian triple is the atomic coordinate-triple contribution to $\sigma_2$ in every ambient
dimension. \Proven

For the full Gaussian--Kronecker row in ambient dimension $n=4$, the tangent space has dimension $3$.
Let $A=S_c$ and $B=S_s$ be the tangent shape-operator pieces induced by $H_c$ and $H_s$, and set
\[
\begin{gathered}
a_1=\operatorname{tr}A,\quad b_1=\operatorname{tr}B,\quad
a_2=\operatorname{tr}A^2,\quad b_2=\operatorname{tr}B^2,\quad
c=\operatorname{tr}(AB),\\
a_3=\operatorname{tr}A^3,\quad b_3=\operatorname{tr}B^3,\quad
d=\operatorname{tr}(A^2B),\quad e=\operatorname{tr}(AB^2).
\end{gathered}
\]
Newton's identities give the explicit four-channel grid
\begin{align*}
K_{3,0}&=\frac{a_1^3-3a_1a_2+2a_3}{6},\\
K_{2,1}&=\frac{a_1^2b_1-2a_1c-b_1a_2+2d}{2},\\
K_{1,2}&=\frac{a_1b_1^2-a_1b_2-2b_1c+2e}{2},\\
K_{0,3}&=\frac{b_1^3-3b_1b_2+2b_3}{6},
\end{align*}
with
\[
K_{GK}=K_{3,0}+K_{2,1}+K_{1,2}+K_{0,3}.
\]
Thus the two mixed orders---two coupling directions against one self direction, and one coupling
direction against two self directions---remain distinct. \Proven

These coefficient identities establish the ambient-$n$ channel account and the explicit $n=4$ row.
They do \emph{not} establish RoleChSpec survival or injectivity for $n\ge4$, classify the general-$n$
mixed-bidegree grid as a representation, or extend an active output chart across a vanishing role
denominator.  Those are separate orbit-separation, mixed-grid, and singular-stratum problems.

\subsection{Curved-ambient extension}
The channel reduction extends to a curved ambient space through the Gauss equation, which reads
schematically
\[
R^\Sigma=\overline R^{\,T}+S\wedge S,
\]
where $R^\Sigma$ is the intrinsic curvature of the role surface, $\overline R^{\,T}$ is the tangential
part of the ambient background curvature, and the extrinsic term is quadratic in the shape operator
$S$. Applying the DBP split to the shape operator, $S=S_c+S_s$, the wedge is bilinear, so
\[
S\wedge S=S_c\wedge S_c+\bigl(S_c\wedge S_s+S_s\wedge S_c\bigr)+S_s\wedge S_s
=R_{cc}+R_{\mathrm{int}}+R_{ss}:
\]
the extrinsic contribution channelises into pure-coupling, interaction, and pure-self parts exactly as
in the flat case. Hence
\[
\text{intrinsic invariant}=\text{ambient input}+\text{extrinsic DBP-channel account}. 
\]
\Proven\ The ambient term $\overline R^{\,T}$ is \emph{not} forced into the same self/coupling split,
and this is a boundary of the construction with a named cause rather than a gap in it: the extrinsic
term is a bilinear function of the shape operator and therefore splits along the role frame, whereas
the ambient curvature is not a bilinear function of the same role data, so no channel split is
available for it unless the application supplies a separate role frame for the ambient geometry
itself. \Proven

%==================================================================
\section{Output-role coupling channels, anisotropy, and faithfulness}
\label{sec:channels-named}

\subsection{The three output-role charts}
For $P=f(D,S)$ with $q_0=1+a^2+b^2$, the three output-role charts carry the channel quadruples
$(\KG,\kc,\ks,\kint)$
\begin{align*}
C_P&=\Bigl(K,\ -\tfrac{B^2}{q_0^2},\ \tfrac{AC}{q_0^2},\ 0\Bigr),\\
C_D&=\Bigl(K,\ -\tfrac{(Ab-aB)^2}{a^2q_0^2},\ \tfrac{A(Ab^2-2abB+a^2C)}{a^2q_0^2},\ 0\Bigr),\\
C_S&=\Bigl(K,\ -\tfrac{(Ca-bB)^2}{b^2q_0^2},\ \tfrac{C(Ca^2-2abB+b^2A)}{b^2q_0^2},\ 0\Bigr),
\end{align*}
with the common Gaussian curvature $K=(AC-B^2)/q_0^2$. The scalar output-role curvature spectrum is
$\operatorname{OutputRoleCurvSpec}(f)=\{C_P,C_D,C_S\}$; retaining ordered input labels gives the
honest sixfold labelled role-jet curvature object. \Symbolic

\subsection{The \texorpdfstring{$n=3$}{n=3} RoleChSpec gauge-obstruction converse}
The preceding chart data may be assembled at orders $1$ and $2$ into the full output-role channel
fingerprint $\operatorname{RoleChSpec}_g(H)$ of a fixed-gradient implicit Hessian: the fingerprint is
the three labelled output-role channel vectors at those two curvature orders.  To distinguish its
gauge-class carrier from the coupling carrier $O$ introduced below, define
\[
\mathcal O_g(H)=(\mathcal O_{12},\mathcal O_{13},\mathcal O_{23}),\qquad
\mathcal O_{ij}=H_{ij}-\frac{g_iH_{jj}}{2g_j}-\frac{g_jH_{ii}}{2g_i},
\]
and $G_g(v)=gv^{\mathsf T}+vg^{\mathsf T}$.

\begin{theorem}[RoleChSpec gauge-obstruction converse in ambient dimension $3$]
\label{thm:rolechspec-converse}
Let $g_1g_2g_3\ne0$ and let $H_1,H_2\in\operatorname{Sym}_3(\mathbb K)$, where
$\operatorname{char}\mathbb K\ne2$.  Then
\[
\operatorname{RoleChSpec}_g(H_1)=\operatorname{RoleChSpec}_g(H_2)
\quad\Longleftrightarrow\quad
\mathcal O_g(H_2-H_1)=0
\quad\Longleftrightarrow\quad
H_2-H_1\in\operatorname{Im}G_g .
\]
Thus the full RoleChSpec is a faithful invariant of the same-gradient gauge quotient
$\operatorname{Sym}_3(\mathbb K)/\operatorname{Im}G_g$ on the regular locus. \Symbolic
\end{theorem}

\begin{proof}
For any $H$, set $v_i=H_{ii}/(2g_i)$.  The representative
$H_\perp=H-G_g(v)$ has zero diagonal and off-diagonal entries
$(\mathcal O_{12},\mathcal O_{13},\mathcal O_{23})$, uniquely; hence
$\ker\mathcal O_g=\operatorname{Im}G_g$.  Direct chart transport shows that RoleChSpec is unchanged by
$G_g(v)$, proving the reverse implication.  Conversely, the six order-$1$ channel components are
linear forms in the three entries of $\mathcal O_g(H)$.  The $3\times3$ minor formed from the three
coupling rows has determinant
\[
\frac{32}{g_1g_2g_3},
\]
which is nonzero on the stated locus when $2$ is invertible.  Equality of the full RoleChSpec therefore
forces equality of the order-$1$ components, and the invertible minor recovers
$\mathcal O_g(H_1)=\mathcal O_g(H_2)$.  Applying linearity to $H_2-H_1$ completes the equivalence.
\end{proof}

This symbolic converse contains the previously frozen $5298$-class exact atlas as a finite
specialisation: the measured survival/collapse biconditional is therefore no longer an open
$n=3$ assertion.  No conclusion is drawn here about RoleChSpec injectivity for $n\ge4$; the rank-$3$
recovery above is dimension-specific.

\subsection{Named coupling channels and the anisotropy scalar}
Define the output-role mixed numerators
\[
\Lambda_P=B,\qquad
\Lambda_D=\frac{Ab-aB}{a},\qquad
\Lambda_S=\frac{Ca-bB}{b},
\]
so that the chart-indexed coupling channels are $\kcrole{\rho}=-\Lambda_\rho^2/q_0^2$ for
$\rho\in\{P,D,S\}$. The coupling carrier is $O=(\kcrole{P},\kcrole{D},\kcrole{S})$. The role anisotropy
scalar is
\[
\Ac=(\kcrole{P}-\kcrole{D})^2+(\kcrole{D}-\kcrole{S})^2+(\kcrole{S}-\kcrole{P})^2,
\]
which vanishes exactly when all three output-role coupling channels agree. \Proven\ Under the
permutation action of the role $S_3$ on $O\in\mathbb R^3$ (permuting the three components),
$\mathbb R^3$ splits as trivial $\oplus$ standard, and a direct computation gives
$\Ac=3\,\lVert P_{\mathbf 2}O\rVert^2$, the squared norm of the standard-irrep part: $\Ac$ is the
intrinsic isotropy defect of the coupling carrier. \Proven\ A product-role excess diagnostic is
$E_P=\kcrole{P}-\tfrac12(\kcrole{D}+\kcrole{S})$, with analogues for $D$ and $S$.

\subsection{Faithfulness of the coupling carrier}
\label{sec:faithful}
The completeness criterion (\S\ref{sec:consequences}) asks whether the carrier separates active role
orbits. The following corollary answers this locally and ties the answer directly to the channels.

\begin{corollary}[faithful coupling carrier]
\label{cor:faithful}
On the regular role locus, the coupling-carrier map of the second jet,
\[
(A,B,C)\longmapsto O=(\kcrole{P},\kcrole{D},\kcrole{S}),
\]
has Jacobian determinant
\[
\det\frac{\partial O}{\partial(A,B,C)}=\frac{8\,\Lambda_P\Lambda_D\Lambda_S}{q_0^{6}} .
\]
Hence the carrier map has rank $3$ wherever $\Lambda_P\Lambda_D\Lambda_S\ne0$; it is then a local
diffeomorphism and therefore separates active role orbits locally, so by the completeness criterion
$O$ is a complete order-$2$ coupling kernel off the channel-isotropy loci. The carrier degenerates
exactly where a coupling numerator vanishes, i.e.\ on the loci where the corresponding channel is
isotropic. \Symbolic
\end{corollary}

\begin{proof}
With $\kcrole{P}=-B^2/q_0^2$, $\kcrole{D}=-(Ab-aB)^2/(a^2q_0^2)$,
$\kcrole{S}=-(Ca-bB)^2/(b^2q_0^2)$ and $q_0$ independent of $(A,B,C)$, the Jacobian has a single
nonzero entry in the first column (from $\kcrole{D}$) and in the third column (from $\kcrole{S}$).
Expanding and simplifying gives the stated determinant; the factorisation is
$8\Lambda_P\Lambda_D\Lambda_S/q_0^6$ on substituting $Ab-aB=a\Lambda_D$, $Ca-bB=b\Lambda_S$,
$B=\Lambda_P$. Full column rank is equivalent to a nonzero determinant, and a nonzero Jacobian gives a
local diffeomorphism, hence local orbit separation; the completeness criterion then applies. The
identity was confirmed in exact symbolic arithmetic.
\end{proof}

Thus the carrier is faithful precisely off the isotropy loci, the same loci on which $\Ac$ records a
collapsing coupling channel; faithfulness and anisotropy are two readings of the channel numerators.

%==================================================================
\section{Gauge transport}
\label{sec:transport}

The reduction $\Sigma$ of \S\ref{sec:channel} has a nontrivial kernel: the zero-sum plane for the
Gaussian case. A channel representative is therefore gauge- and frame-relative, free to move inside
the fibre over the fixed invariant. Concretely, while the graph-channel formula sets $\kint=0$ in any
graph gauge, transport between graph sections of the same relation, or passage to a non-graph
implicit gauge, moves weight between $\kc$, $\kint$ and $\ks$ along $\ker\Sigma$, with $\KG=\sigma_2$
held fixed. The interaction channel is therefore live precisely when the working gauge is not a graph
gauge; the keystone of Appendix~\ref{app:examples} exhibits an implicit relation whose entire
curvature is carried by $\kint$. \Symbolic

%==================================================================
\section{Singular strata}
\label{sec:singular}

The active maps are defined only where their denominators are nonzero. Where a required denominator
vanishes the corresponding output chart does not exist, and the correct object is a typed singular
stratum, not a coerced scalar. Typical cases:
\begin{enumerate}[itemsep=1pt,topsep=2pt]
\item a coordinate cannot be solved locally as output;
\item a projective normal coordinate vanishes;
\item a curvature denominator degenerates;
\item a requested channel representative is not defined in the chosen frame.
\end{enumerate}
A singular stratum may encode a symmetry, a boundary, a chart failure, or a refusal condition; in all
cases it is part of the role-jet structure rather than an exception to be removed. \Proven

%==================================================================
\section{Scope and summary}
\label{sec:scope}

The theorem is local and finite-order: it describes order-$r$ jets near a regular point, and global
statements require additional gluing data across charts. It is frame-relative: the channel split is
made in the distinguished coordinate frame, so the reduced invariant may be intrinsic while a channel
account is basis-, chart-, or gauge-relative. The theorem does not assert that every useful scalar is
invariant; it gives a test, namely that a scalar is a DBP invariant exactly when it is constant on the
active role orbits of the relevant datum.

\medskip
\noindent In compact form:
\[
\boxed{\;
\begin{gathered}
\text{DBP invariant data is carried by the active role-jet orbit;}\\
\text{DBP channel data is the fibre over its invariant reduction.}
\end{gathered}\;}
\]

%==================================================================
\appendix
\section{Higher-order transposition}
\label{app:higher}

The active transposition $t=(D\;P)$ on the order-$r$ jet is obtained from $D=q(P,S)$ by a homogeneous
inverse recurrence.  Translate the base point to the origin and write
\[
f(D,S)=aD+bS+\sum_{d=2}^{r}f_d(D,S),
\qquad
q(P,S)=a^{-1}P-\frac baS+\sum_{d=2}^{r}q_d(P,S),
\]
where $f_d$ and $q_d$ are homogeneous of total degree $d$ and $a=f_D\ne0$.  Suppose
$q_2,\ldots,q_{d-1}$ have been chosen and put
\[
q_{<d}=a^{-1}P-\frac baS+\sum_{j=2}^{d-1}q_j,
\qquad
R_d=[\,f(q_{<d}(P,S),S)-P\,]_d .
\]
Here $[\,\cdot\,]_d$ denotes the homogeneous degree-$d$ part.  Replacing $q_{<d}$ by
$q_{<d}+q_d$ changes that degree-$d$ residual by exactly $a q_d$: inserting $q_d$ into any nonlinear
$f_j$, $j\ge2$, has degree at least $d+1$.  The unique cancelling choice is therefore
\[
q_d=-a^{-1}R_d.
\]
Induction on $d=2,\ldots,r$ constructs the unique inverse jet and shows
\[
q_d\in\mathbb Q[\,\text{coefficients of }f_1,\ldots,f_d,\ a^{-1}],
\]
so every order-$r$ component is rational in the original jet, with only powers of the first-order role
denominator.  In particular $q_d$ depends on no input coefficient above degree $d$.  If
$\tau_{N,r}$ denotes truncation from order $N$ to order $r$, this triangular dependence gives
\[
\tau_{N,r}(t f)=t\tau_{N,r}(f).
\]
The input swap $s$ preserves total degree and likewise commutes with truncation.  An induction on the
length of a word $w\in\langle s,t\rangle$ then yields
\[
\tau_{N,r}(w f)=w\tau_{N,r}(f)
\qquad(N\ge r\ge2),
\]
which writes out the all-finite-order active-role induction.  The first two orders are the formulas of
\S\ref{sec:transpose}; the recurrence covers every finite order on the same regular chart locus.
\Proven

%==================================================================
\section{Worked examples}
\label{app:examples}

\subsection{Power-law charts: value-only role spectra}
For a power-law coupling $P=D^mS^n$ pass to log coordinates, $\log P=m\log D+n\log S$, with exponent
normal $\ell=(-m,-n,1)$. Role transposition permutes the entries of this projective normal, giving the
exponent-pair orbit
\[
\mathcal E(m,n)=\Bigl\{(m,n),(n,m),\bigl(\tfrac1m,-\tfrac nm\bigr),\bigl(-\tfrac nm,\tfrac1m\bigr),
\bigl(-\tfrac mn,\tfrac1n\bigr),\bigl(\tfrac1n,-\tfrac mn\bigr)\Bigr\}
\]
and the role-degree spectrum
\[
\operatorname{DegSpec}(m,n)=\Bigl\{m+n,\ m+n,\ \tfrac{1-n}{m},\ \tfrac{1-n}{m},\ \tfrac{1-m}{n},\
\tfrac{1-m}{n}\Bigr\}.
\]
For ordinary multiplication $P=DS$, $(m,n)=(1,1)$ gives $\operatorname{DegSpec}(1,1)=\{2,2,0,0,0,0\}$:
the forward chart has degree $2$ while inverse charts $D=P/S$ have degree $0$ in log-role form. For
the cubic mixed power $P=D^2S$, $(m,n)=(2,1)$ gives
$\operatorname{DegSpec}(2,1)=\{3,3,0,0,-1,-1\}$, distinguishing it from ordinary multiplication though
both are monomial. Scale-free exponent invariants are
\[
I_1^\ell=\frac{(1-m-n)(mn-m-n)}{mn},\qquad
I_2^\ell=\frac{(1-m-n)^3}{mn},
\]
with exponent discriminant $\Delta_\ell=(m-n)^2(m+1)^2(n+1)^2$. \Symbolic\ At
multiplication the coupling channel is flat, so $\Ac=0$; the keystone below is nowhere flat.

\subsection{The keystone: interaction-only curvature}
Consider the implicit relation
\[
F=x_1^2+x_1x_2+x_3^2-3\qquad\text{at}\qquad p=(1,1,1).
\]
Then $g=(3,1,2)$, $q=14$, and
\[
H=\begin{pmatrix}2&1&0\\1&0&0\\0&0&2\end{pmatrix}.
\]
The channel-density polynomial of \S\ref{sec:channel} gives
\[
\kc=-\tfrac1{49},\qquad \ks=\tfrac1{49},\qquad \kint=-\tfrac3{49},
\qquad\text{so}\qquad \KG=\kc+\kint+\ks=-\tfrac3{49}.
\]
The coupling and self channels cancel exactly and the surviving curvature is interaction curvature.
Any non-negative scalar proxy, or any single-channel account that omits $\kint$, cannot represent this
object; the keystone is the discriminating witness that the three-channel decomposition is not
reducible to a single number. \Symbolic\ (Channels confirmed in exact arithmetic against the direct
Gauss--Kronecker value $\KG=-3/49$.)

\subsection{Gauge transport on the keystone}
A defining-function gauge $\widetilde F=\mu F$ with $\mu(p)=1$ and $a=\nabla\log\mu(p)$ fixes the
gradient and shifts the Hessian by
\[
\widetilde g=g,\qquad \widetilde H=H+ga^{\mathsf T}+ag^{\mathsf T}.
\]
Writing the gauge vector as $a=(u,v,w)$ on the keystone surface above, the channels transport by the
exact rational laws
\begin{align*}
\kc(u,v,w)&=\tfrac1{49}\bigl(12uv+6uw+18vw+6w-1\bigr),\\
\ks(u,v,w)&=\tfrac1{49}\bigl(12uv+6uw+3u+18vw+13v+2w+1\bigr),\\
\kint(u,v,w)&=-\tfrac1{49}\bigl(24uv+12uw+3u+36vw+13v+8w+3\bigr),
\end{align*}
and for every rational gauge vector
\[
\kc(u,v,w)+\ks(u,v,w)+\kint(u,v,w)=-\tfrac3{49}.
\]
The channel triple therefore moves inside the affine zero-sum plane
$\Delta\kc+\Delta\ks+\Delta\kint=0$ while the Gaussian curvature is held fixed, concretely realising
the kernel of the reduction $\Sigma$ (\S\ref{sec:channel}) and the gauge transport of
\S\ref{sec:transport}. The pure coupling channel is pinned, $\kc=-\tfrac1{49}$, exactly on the locus
\[
w(u+3v+1)+2uv=0,
\]
so channel pinning is a structural condition on the gauge, not an accident: a generic gauge moves
every channel, a special gauge fixes a chosen one. \Symbolic

%==================================================================
\section{Application patterns and working checklist}
\label{app:patterns}

The calculus is applied by supplying a role assignment and a relation and then reading the orbit and
channel account at the order the question requires.
\begin{itemize}[itemsep=2pt,topsep=2pt]
\item \textbf{Generic relation.} Assign $D$ (directive, policy, or design pressure), $S$ (substrate,
state, or resource), $P$ (product, output, or reading), and write $P=f(D,S)$ or $F(D,S,P)=0$. The role
orbit records what survives re-solving from each role; the channel account records how the invariant
is distributed in the chosen frame.
\item \textbf{Sensitivity.} If only first order matters, take $(a,b)=(f_D,f_S)$ and its active orbit;
projective role invariants diagnose role symmetry, stabilisers, and inverse-role boundaries.
\item \textbf{Curvature.} If second-order interaction matters, take $(a,b,A,B,C)$ and read the
output-role spectra $C_P,C_D,C_S$, with $\KG$ the reduced invariant and the channel triples the
explanatory representatives.
\item \textbf{Route/account.} If several routes reach the same reading, treat route choice as a
role-structured relation; the invariant quotient is route-independent, and a loop holonomy is a
coupling diagnostic that vanishes when routes combine additively and is nonzero when they genuinely
interact.
\item \textbf{Monomial scaling.} For $P=D^mS^n$, use the exponent orbit and degree spectrum to detect
when a monomial that is simple in one role becomes inverse, neutral, or singular in another
(Appendix~\ref{app:examples}).
\end{itemize}

\noindent A working checklist for a new application:
\begin{center}
\small
\begin{tabular}{r l l}
\hline
& Object & Question\\
\hline
1 & Roles $(D,S,P)$ & What are the directive, substrate, product roles?\\
2 & Relation & Is the local law $P=f(D,S)$ or $F(D,S,P)=0$?\\
3 & Regular locus & Which role denominators must be nonzero?\\
4 & Jet order $r$ & Sensitivity, curvature, or higher holonomy?\\
5 & Active role orbit & What are the transformed role charts?\\
6 & Candidate invariant & Does it factor through the active role quotient?\\
7 & Channel account & Which explanatory channels are retained?\\
8 & Reduction map & What is the invariant quotient of the account?\\
9 & Kernel / fibre & What moves under chart or gauge while the invariant stays fixed?\\
10 & Singular strata & What must refuse, rechart, or be typed at a boundary?\\
\hline
\end{tabular}
\end{center}

%==================================================================
\section*{Provenance of claims}
\addcontentsline{toc}{section}{Provenance of claims}

The $S_3$ relations, the first- and second-order transposition formulas, and the homogeneous
all-finite-order inverse recurrence are explicit rational identities (\Proven, independently replayed
coefficient-by-coefficient). Theorem~\ref{thm:orbit}, the completeness criterion,
the exact rational closure, the channel-reduction exact sequence, the ambient-$n$ bordered-minor
channel law, the explicit $n=4$ Gaussian--Kronecker grid, the Gaussian decomposition
$\KG=\kc+\kint+\ks$, the graph-channel formula, the bilinear curved-ambient split
$S\wedge S=R_{cc}+R_{\mathrm{int}}+R_{ss}$ with its named ambient boundary, and the definitions and
identities of the named channels and $\Ac=3\lVert P_{\mathbf 2}O\rVert^2$ are \Proven.
Theorem~\ref{thm:rolechspec-converse}, Corollary~\ref{cor:faithful} (the carrier Jacobian determinant
$8\Lambda_P\Lambda_D\Lambda_S/q_0^6$ and rank $3$), the keystone channel values, and the keystone
gauge-transport laws were established and confirmed in exact symbolic arithmetic, \Symbolic. The
RoleChSpec converse is confined to ambient dimension $3$ and makes no $n\ge4$ injectivity claim. The
construction is self-contained: no external
named theorem is load-bearing, and every transported quantity is exact-$\mathbb Q$ on the regular
locus by \S\ref{sec:consequences}.

\end{document}
